An MDPI-structured research article must state where its data and code
live and how a reader can reproduce the reported results. This chapter
provides that statement plus the mapping from each numeric claim in the
paper to the raw per-iteration CSV that supports it.

\begin{center}\rule{0.5\linewidth}{0.5pt}\end{center}

\section{1. Data availability
statement}\label{data-availability-statement}

All data presented in this study are openly archived alongside the
source code of the measurement framework in this repository. The raw
per-iteration CSVs for every experiment are under
\href{../experiments/}{\texttt{../experiments/}}, mirrored byte-for-byte
from the device's LittleFS \texttt{/diags/} tree at the time of capture.
Statistical summaries in the paper tables are derived offline by the
generator script
\href{../experiments/_build_appendix.py}{\texttt{../experiments/\_build\_appendix.py}},
which operates solely on the checked-in CSVs and reproduces every
appendix number without access to the hardware.

No proprietary, subject-private, or licence-restricted data is involved.

\begin{center}\rule{0.5\linewidth}{0.5pt}\end{center}

\section{2. Repository layout}\label{repository-layout}

At repository root, the subtrees that together constitute the reported
artifact are:

{\def\LTcaptype{none} % do not increment counter
\begin{longtable}[]{@{}
  >{\raggedright\arraybackslash}p{(\linewidth - 2\tabcolsep) * \real{0.5000}}
  >{\raggedright\arraybackslash}p{(\linewidth - 2\tabcolsep) * \real{0.5000}}@{}}
\toprule\noalign{}
\begin{minipage}[b]{\linewidth}\raggedright
Path
\end{minipage} & \begin{minipage}[b]{\linewidth}\raggedright
Purpose
\end{minipage} \\
\midrule\noalign{}
\endhead
\bottomrule\noalign{}
\endlastfoot
\texttt{examples/sentai\_runtime/} & Firmware application source
(\texttt{sentai\_runtime.cc}, \texttt{modsentai\_*.c},
\texttt{detection\_task.cc}, \texttt{sentai\_httpd.cc},
\texttt{sentai\_lfs\_task.cc}, \ldots) \\
\texttt{examples/sentai\_runtime/diag/} & MicroPython diagnostic package
--- experiments E1-E18 \\
\texttt{examples/sentai\_runtime/diag/\_host\_upload\_repl.py} &
Host-side REPL uploader used to push \texttt{diag/*.py} to the board \\
\texttt{examples/sentai\_runtime/diag/drivers/} & Host-side REPL drivers
(\texttt{\_e18\_post\_refactor.py}, \ldots) that invoke experiments
end-to-end \\
\texttt{examples/sentai\_runtime/experiments/} & The archived session
data presented in this paper \\
\texttt{examples/sentai\_runtime/paper/} & Paper source (this file and
siblings) \\
\texttt{examples/sentai\_runtime/agent/} & Operating documentation
(embeded.md discipline rules, agent.md handoff guide) \\
\texttt{libs/base/gpio.cc}, \texttt{libs/camera/},
\texttt{libs/base/filesystem.cc}, \ldots{} & SentAI platform libraries
modified as part of this work \\
\texttt{third\_party/nxp/rt1176-sdk/…} & Vendor SDK, used as-is except
for documented table entries \\
\texttt{third\_party/micropython/} & MicroPython embed port, used
as-is \\
\texttt{scripts/flashtool.py} & NXP flashing tool used to program the
board \\
\end{longtable}
}

A reader who wishes to reproduce the measurements in this paper needs
the three subtrees \texttt{examples/sentai\_runtime/*} plus the
cross-cutting library changes in \texttt{libs/camera/cam\_mux.h} and
\texttt{libs/base/gpio.cc}, all of which are under the repository's
open-source licence.

\begin{center}\rule{0.5\linewidth}{0.5pt}\end{center}

\section{3. Firmware build
identification}\label{firmware-build-identification}

The firmware image presented as the ``Fix B post-review'' configuration
--- responsible for every numeric claim in Tables 3, 4, 6 and 7 of
\url{evaluation.md} --- is:

{\def\LTcaptype{none} % do not increment counter
\begin{longtable}[]{@{}
  >{\raggedright\arraybackslash}p{(\linewidth - 2\tabcolsep) * \real{0.5000}}
  >{\raggedright\arraybackslash}p{(\linewidth - 2\tabcolsep) * \real{0.5000}}@{}}
\toprule\noalign{}
\begin{minipage}[b]{\linewidth}\raggedright
Field
\end{minipage} & \begin{minipage}[b]{\linewidth}\raggedright
Value
\end{minipage} \\
\midrule\noalign{}
\endhead
\bottomrule\noalign{}
\endlastfoot
Build number & 640 \\
Build timestamp & stored in \texttt{build\_version.h} and echoed at
every boot: \texttt{SentAI\ build\ \#640\ (2026-04-20\ …)} \\
Git branch & \texttt{feature/ov5640-camera-support} \\
Full toolchain & CMake 3.x + Ninja, \texttt{arm-none-eabi-gcc} 10.3.x,
NXP MCUXpresso SDK vendored under
\texttt{third\_party/nxp/rt1176-sdk/} \\
Reconstruction &
\texttt{cmake\ -S\ .\ -B\ build\ \&\&\ cmake\ -\/-build\ build\ -\/-target\ sentai\_runtime} \\
Flash & \texttt{python3\ scripts/flashtool.py\ -e\ sentai\_runtime}
(persistent; omit \texttt{-\/-ram}) \\
\end{longtable}
}

Earlier builds (\#622 for the first eDMA measurement; \#632 for Fix A;
\#635 for Fix B pre-refactor) are reachable by checking out earlier
commits of the same branch and rebuilding. The cross-session
reproducibility table (Table 4 of \url{evaluation.md}) documents the
three builds that together demonstrate the refactor is
performance-neutral.

\begin{center}\rule{0.5\linewidth}{0.5pt}\end{center}

\section{4. Hardware requirements}\label{hardware-requirements}

{\def\LTcaptype{none} % do not increment counter
\begin{longtable}[]{@{}
  >{\raggedright\arraybackslash}p{(\linewidth - 2\tabcolsep) * \real{0.5000}}
  >{\raggedright\arraybackslash}p{(\linewidth - 2\tabcolsep) * \real{0.5000}}@{}}
\toprule\noalign{}
\begin{minipage}[b]{\linewidth}\raggedright
Requirement
\end{minipage} & \begin{minipage}[b]{\linewidth}\raggedright
Minimum
\end{minipage} \\
\midrule\noalign{}
\endhead
\bottomrule\noalign{}
\endlastfoot
MCU board & Coral Dev Board Micro with SentAI daughter-board v1.0 \\
USB cable + host & one USB-C cable to a Linux host (CDC-ACM + CDC-NCM
required; most host USB stacks supply these out of the box) \\
Scene & static scene with either no target class or a fixed target
count; the paper reports a zero-target scene --- see
\url{experimental_setup.md} §5 \\
Bench time & ≈ 5 s per experiment sweep; ≈ 30 s per session including
begin/end/snapshots; the entire reported E18 benchmark is ≈ 15 s
on-device \\
\end{longtable}
}

No special test equipment (oscilloscope, logic analyser, thermal
chamber) is required to reproduce any number in this paper. That is
intentional: the measurement framework captures every stage it needs
internally and emits CSVs.

\begin{center}\rule{0.5\linewidth}{0.5pt}\end{center}

\section{5. Reproducing a single session from
scratch}\label{reproducing-a-single-session-from-scratch}

Assuming the repository is checked out at the branch and the board is
connected:

\begin{verbatim}
# 1. Build + flash
cd /path/to/coralmicro
cmake -S . -B build
cmake --build build --target sentai_runtime
python3 scripts/flashtool.py -e sentai_runtime

# 2. Wait for the board to boot (up to ~15 s first time after flash)
for i in $(seq 1 20); do
  curl -s -m 2 -o /dev/null http://10.0.0.1/ && break
  sleep 1
done

# 3. Push the latest diag/ package
cd examples/sentai_runtime
python3 diag/_host_upload_repl.py \
    --file e_pipeline.py --file _util.py \
    --file _session.py --file __init__.py

# 4. Execute the reported head-to-tail benchmark
python3 diag/drivers/_e18_post_refactor.py

# 5. Fetch the session folder (session id is the next free sNNN on
#    this board's counter; inspect /api/ls/diags/ to find it)
curl -s http://10.0.0.1/api/ls/diags/ | python3 -m json.tool | less

# 6. Pull the three CSVs
for f in 001_e18_A_fixed_cam0.csv 002_e18_B_fixed_cam1.csv \
         003_e18_C_alt_cam0_cam1.csv; do
    curl -s http://10.0.0.1/api/raw/diags/sNNN_e18_post_refactor/$f > $f
done
\end{verbatim}

The numbers from step 6 should agree with Table 3 of \url{evaluation.md}
within the ≤ 1 ms cross-session noise documented in Table 4.

\begin{center}\rule{0.5\linewidth}{0.5pt}\end{center}

\section{6. Claim-to-file map (the ``which CSV supports which table''
reference)}\label{claim-to-file-map-the-which-csv-supports-which-table-reference}

The table below is the single authoritative map from every numeric claim
in this paper to its on-disk evidence. Rows are ordered by where the
claim appears in the paper.

\textbf{Table 1.} Paper-level claim-to-evidence index.

{\def\LTcaptype{none} % do not increment counter
\begin{longtable}[]{@{}
  >{\raggedright\arraybackslash}p{(\linewidth - 6\tabcolsep) * \real{0.2500}}
  >{\raggedright\arraybackslash}p{(\linewidth - 6\tabcolsep) * \real{0.2500}}
  >{\raggedright\arraybackslash}p{(\linewidth - 6\tabcolsep) * \real{0.2500}}
  >{\raggedright\arraybackslash}p{(\linewidth - 6\tabcolsep) * \real{0.2500}}@{}}
\toprule\noalign{}
\begin{minipage}[b]{\linewidth}\raggedright
Claim
\end{minipage} & \begin{minipage}[b]{\linewidth}\raggedright
Where it appears
\end{minipage} & \begin{minipage}[b]{\linewidth}\raggedright
Evidence (folder + file)
\end{minipage} & \begin{minipage}[b]{\linewidth}\raggedright
Generator
\end{minipage} \\
\midrule\noalign{}
\endhead
\bottomrule\noalign{}
\endlastfoot
CPU memcpy ≈ 24 ms baseline & \url{memcpy.md} §``Baseline'',
\url{evaluation.md} Table 2 & \href{memcpy.md}{memcpy.md CSV block}
(10-run table is reproduced inline); raw per-iteration CSVs captured by
\href{../_e15_ab.py}{\texttt{../\_e15\_ab.py}} during an A/B session,
archived inline in the paper chapter & inline \\
eDMA memcpy ≈ 14.6 ms & \url{memcpy.md} §``Optimised'',
\url{evaluation.md} Table 2 & same A/B session as above & inline \\
15.47 FPS post-eDMA sustained, 512×512 & \url{memcpy.md} §``Per-run
results'' optimised block; \url{evaluation.md} Table 1 &
\href{../experiments/s032_e15_x20/}{\texttt{../experiments/s032\_e15\_x20/}}
(20 × 20 frames, post-eDMA) &
\href{../experiments/_build_appendix.py}{\texttt{\_build\_appendix.py}}
§Appendix C \\
13.39 FPS pre-eDMA & \url{memcpy.md} §``Baseline'' & same A/B session,
baseline block & inline \\
E14 baseline 6.4 FPS (80-class model) & \url{evaluation.md} §8 ``Each
optimisation exposes the next''; Appendix B &
\href{../experiments/s021_e14_x20/}{\texttt{../experiments/s021\_e14\_x20/}},
\href{../experiments/s022_e14_x20/}{s022},
\href{../experiments/s023_e14_x20/}{s023},
\href{../experiments/s025_e14_x20/}{s025},
\href{../experiments/s031_e14_x20/}{s031} &
\href{../experiments/_build_appendix.py}{\texttt{\_build\_appendix.py}}
§Appendix B \\
Pre-Fix A alternating FPS 4.7, asymmetry +64 ms & \url{cam_switch.md}
§``E15 vs E16 --- head-to-head''; \url{evaluation.md} Table 7 &
\href{../experiments/s034_e15_vs_e16_x40/002_e16_camswitch_cam0_cam1.csv}{\texttt{../experiments/s034\_e15\_vs\_e16\_x40/002\_e16\_camswitch\_cam0\_cam1.csv}}
&
\href{../experiments/_build_appendix.py}{\texttt{\_build\_appendix.py}}
§Appendix D \\
Fix A null result (asymmetry unchanged) & \url{cam_switch.md} §``Fix
A''; \url{evaluation.md} Table 7 &
\href{../experiments/s035_e15_vs_e16_x40/002_e16_camswitch_cam0_cam1.csv}{\texttt{../experiments/s035\_e15\_vs\_e16\_x40/002\_e16\_camswitch\_cam0\_cam1.csv}}
&
\href{../experiments/_build_appendix.py}{\texttt{\_build\_appendix.py}}
§Appendix D \\
Fix B alternating FPS 6.87, asymmetry ≤ 1 ms (E16 alone) &
\url{cam_switch.md} §``Fix B --- flip-on-EOF'' &
\href{../experiments/s042_e16_eof_30fps_x40/001_e16_camswitch_cam0_cam1.csv}{\texttt{../experiments/s042\_e16\_eof\_30fps\_x40/001\_e16\_camswitch\_cam0\_cam1.csv}}
&
\href{../experiments/_build_appendix.py}{\texttt{\_build\_appendix.py}}
§Appendix E \\
E18 head-to-tail, pre-refactor (s043) & \url{cam_switch.md}
§``Head-to-tail benchmark'' &
\href{../experiments/s043_e18_headtail_drain2/}{\texttt{../experiments/s043\_e18\_headtail\_drain2/\{001,002,003\}\_e18\_*.csv}}
&
\href{../experiments/_build_appendix.py}{\texttt{\_build\_appendix.py}}
§Appendix G \\
E18 head-to-tail, post-refactor (s045) & \url{cam_switch.md} §``Session
\texttt{s045…}''; \url{evaluation.md} Table 3 &
\href{../experiments/s045_e18_post_refactor/}{\texttt{../experiments/s045\_e18\_post\_refactor/\{001,002,003\}\_e18\_*.csv}}
&
\href{../experiments/_build_appendix.py}{\texttt{\_build\_appendix.py}}
§Appendix G \\
Cross-session reproducibility ≤ 1 ms & \url{evaluation.md} Table 4 &
\href{../experiments/s043_e18_headtail_drain2/}{\texttt{../experiments/s043\_e18\_headtail\_drain2/}},
\href{../experiments/s044_e18_headtail_drain2/}{\texttt{s044}},
\href{../experiments/s045_e18_post_refactor/}{\texttt{s045}} &
derivation in \url{statistical_notes.md} \\
Pre-fix mid-buffer seam (visual) & \url{cam_switch.md} §``Visual
evidence''; \url{evaluation.md} Table 5 &
\href{../experiments/s038_e17_drain_ab/e17_t1_frames/}{\texttt{../experiments/s038\_e17\_drain\_ab/e17\_t1\_frames/\{002\_cam0\_133ms,003\_cam1\_202ms\}.jpg}}
& visual inspection \\
Post-fix seam-free at drain=2 (visual) & \url{cam_switch.md} §``Visual
evidence''; \url{evaluation.md} Table 5 &
\href{../experiments/s041_e17_eof_check/e17_t2_frames/}{\texttt{../experiments/s041\_e17\_eof\_check/e17\_t2\_frames/}}
& visual inspection \\
drain=1 residual artefacts & \url{cam_switch.md} §``Known limitations'';
\url{threats_to_validity.md} §2.5 &
\href{../experiments/s041_e17_eof_check/e17_t1_frames/}{\texttt{../experiments/s041\_e17\_eof\_check/e17\_t1\_frames/}}
& visual inspection \\
Fault counters zero on nominal run & \url{evaluation.md} Table 6;
\url{cam_switch.md} §``Shipping runtime surface'' &
\texttt{sentai.diag.cam\_stats()} output at end of sessions
\texttt{s043}, \texttt{s044}, \texttt{s045} &
\href{../modsentai_diag.c}{modsentai\_diag.c:mod\_sentai\_diag\_cam\_stats} \\
\end{longtable}
}

\begin{center}\rule{0.5\linewidth}{0.5pt}\end{center}

\section{7. Data integrity check}\label{data-integrity-check}

A reviewer can verify the integrity of the archived CSVs against paper
tables with three commands:

\begin{verbatim}
# 1. All appendix tables re-derived from raw CSVs (no board access)
cd examples/sentai_runtime
python3 experiments/_build_appendix.py  # overwrites experiments/_appendix section of README

# 2. The specific E18 head-to-tail claim
python3 - <<'EOF'
import csv, statistics
def mean(csvpath, col):
    rows = list(csv.DictReader(open(csvpath).readlines()))[1:]  # drop warmup
    return statistics.mean(int(r[col]) for r in rows)
for sess in ('s043_e18_headtail_drain2', 's044_e18_headtail_drain2', 's045_e18_post_refactor'):
    for (label, f) in (('A', '001_e18_A_fixed_cam0.csv'),
                       ('B', '002_e18_B_fixed_cam1.csv'),
                       ('C', '003_e18_C_alt_cam0_cam1.csv')):
        m = mean('experiments/%s/%s' % (sess, f), 'total_frame_ms')
        print('%s %s mean=%.1f fps=%.2f' % (sess, label, m, 1000/m))
EOF

# 3. Visual evidence: open the two representative frames and compare
xdg-open experiments/s038_e17_drain_ab/e17_t1_frames/002_cam0_133ms.jpg
xdg-open experiments/s041_e17_eof_check/e17_t1_frames/003_cam1_201ms.jpg
\end{verbatim}

\begin{center}\rule{0.5\linewidth}{0.5pt}\end{center}

\section{8. Author contributions, funding, conflicts (MDPI
boilerplate)}\label{author-contributions-funding-conflicts-mdpi-boilerplate}

\emph{(Placeholders. These will be filled by the authors at submission.
Included here so the paper source already has the MDPI-required
declarations in a canonical location, rather than being scattered across
templates.)}

\begin{itemize}
\tightlist
\item
  \textbf{Author Contributions}: Conceptualization, \emph{TBD};
  methodology, \emph{TBD}; software, \emph{TBD}; validation, \emph{TBD};
  formal analysis, \emph{TBD}; investigation, \emph{TBD}; resources,
  \emph{TBD}; data curation, \emph{TBD}; writing --- original draft
  preparation, \emph{TBD}; writing --- review and editing, \emph{TBD};
  visualization, \emph{TBD}; supervision, \emph{TBD}; project
  administration, \emph{TBD}; funding acquisition, \emph{TBD}. All
  authors have read and agreed to the published version of the
  manuscript.
\item
  \textbf{Funding}: \emph{This research received no external funding.}
  (or list the specific grant).
\item
  \textbf{Institutional Review Board Statement}: \emph{Not applicable.}
  (This study involves no human subjects, no animal subjects, and no
  identifying data --- it is an embedded-systems measurement study on a
  development board.)
\item
  \textbf{Informed Consent Statement}: \emph{Not applicable.}
\item
  \textbf{Data Availability Statement}: All data presented in this study
  are available in the repository accompanying this paper under
  \texttt{examples/sentai\_runtime/experiments/}, along with the
  generator script that reproduces every summary statistic from the raw
  per-iteration CSVs. See §1 above for details.
\item
  \textbf{Acknowledgments}: \emph{TBD} --- list any non-funding support
  (infrastructure, advice, early reviewers).
\item
  \textbf{Conflicts of Interest}: \emph{The authors declare no conflicts
  of interest.} (or as applicable).
\end{itemize}

\begin{center}\rule{0.5\linewidth}{0.5pt}\end{center}

\section{9. Supplementary materials}\label{supplementary-materials}

Following MDPI's supplementary-materials convention, these artefacts
accompany the main paper but are not required to read the narrative:

\begin{itemize}
\tightlist
\item
  \texttt{experiments/README.md} --- per-session narrative index +
  appendices A--G with per-experiment tables derived from every CSV in
  this archive.
\item
  \texttt{experiments/methodology.md} --- full measurement protocol
  (warm-up, drop-first-sample convention, reproducibility rules, noise
  budget).
\item
  \texttt{experiments/\_build\_appendix.py} --- offline appendix
  generator; runs on the CSVs in this archive, produces the appendix
  markdown with no board access.
\item
  \texttt{agent/agent.md} --- operational handoff guide (how to pick up
  this project, how to drive the REPL, how to regenerate QSTRs, how to
  warm-reset a stuck REPL).
\item
  \texttt{agent/embeded.md} --- NASA/JPL-style coding-discipline rules
  used during the firmware refactor phase.
\end{itemize}

All of the above are in-repository, tracked by git, and version- aligned
with the firmware that produced the reported measurements.
